%----------------------------------------------------------------
% low-level LaTeX programming
%----------------------------------------------------------------

\newcommand\hmmax{0}
\newcommand\bmmax{0}

% e-TeX tools for LaTeX
% https://ctan.org/pkg/etoolbox
\usepackage{etoolbox}

% macro patching commands
% https://ctan.org/pkg/xpatch
\usepackage{xpatch}

% a generic document command parser
% https://ctan.org/pkg/xparse
\usepackage{xparse}

% a new interface for environments in LaTeX
% https://ctan.org/pkg/environ
\usepackage{environ}

% driver-independent color extensions for LaTeX and pdfLaTeX
% https://ctan.org/pkg/xcolor
\usepackage{xcolor}

% print lengths using specified units
% https://ctan.org/pkg/printlen
\usepackage{printlen}

% macro package for creating graphics
% https://ctan.org/pkg/tikz
\usepackage{tikz}
\usetikzlibrary{positioning,shapes,arrows.meta,calc}

% control rendering parameters
% https://ctan.org/pkg/pdfrender
\usepackage{pdfrender}

%----------------------------------------------------------------
% language and fonts
%----------------------------------------------------------------

% select font encodings
% https://www.ctan.org/pkg/fontenc
\usepackage[T1]{fontenc}

% accept different input encodings
% https://www.ctan.org/pkg/inputenc
\usepackage[utf8]{inputenc}

% manages culturally-determined typographical (and other) rules for a wide range of languages
% https://www.ctan.org/pkg/babel
\usepackage[english]{babel}

% use Times as default text font, and provide maths support
% https://ctan.org/pkg/mathptmx
% \usepackage{mathptmx}
\usepackage{times}
% \usepackage{newtx}

% Font Awesome 5 with LaTeX support
% https://ctan.org/pkg/fontawesome5
\usepackage{fontawesome5}

% LaTeX support for the Text Companion fonts
% https://ctan.org/pkg/textcomp
% \usepackage{textcomp}

% access to PostScript standard Symbol and Dingbats fonts
% https://ctan.org/pkg/pifont
% \usepackage{pifont}

%----------------------------------------------------------------
% text formatting
%----------------------------------------------------------------

% subliminal refinements towards typographical perfection
% https://ctan.org/pkg/microtype
\usepackage[stretch=10]{microtype}

% margin adjustment and detection of odd/even pages
% https://ctan.org/pkg/changepage
\usepackage{changepage}

% set space between lines
% https://ctan.org/pkg/setspace
% \usepackage{setspace}

% alternative versions of “ragged”-type commands
% https://ctan.org/pkg/ragged2e
\usepackage{ragged2e}

% package for underlining
% https://ctan.org/pkg/ulem
\usepackage[normalem]{ulem}

% print a coloured contour around text
% https://ctan.org/pkg/contour
\usepackage{contour}

%----------------------------------------------------------------
% math formatting
%----------------------------------------------------------------

% https://tex.stackexchange.com/questions/32100
% \usepackage{amsmath} % mathtools below
\usepackage{amssymb}
\usepackage{amsthm}

% mathematical tools to use with amsmath
% https://ctan.org/pkg/mathtools
\usepackage{mathtools}

% access bold symbols in maths mode
% https://ctan.org/pkg/bm
\usepackage{bm}

% improve the typesetting of mathematical matrices with PGF
% https://ctan.org/pkg/nicematrix
\usepackage{nicematrix}

% multiple mathematical accents
% https://ctan.org/pkg/accents
\usepackage{accents}

% typeset in-line fractions in a "nice" way
% https://ctan.org/pkg/nicefrac
\usepackage{nicefrac}

% place lines through maths formulae
% https://ctan.org/pkg/cancel
\usepackage{cancel}

% commands to produce dots in math that respect font size
% https://ctan.org/pkg/mathdots
\usepackage{mathdots}

%----------------------------------------------------------------
% float and similar object formatting
%----------------------------------------------------------------

% enhanced support for graphics
% https://ctan.org/pkg/graphicx
\usepackage{graphicx}

% improved interface for floating objects
% https://ctan.org/pkg/float
% \usepackage{float}

% graphics package-alike macros for “general” boxes
% https://ctan.org/pkg/adjustbox
\usepackage{adjustbox}

% combine LaTeX commands over included graphics
% https://ctan.org/pkg/overpic
\usepackage[percent]{overpic}

% support for sub-captions
% https://ctan.org/pkg/subcaption
% \usepackage{subcaption}

% Publication quality tables in LaTeX
% https://ctan.org/pkg/booktabs
\usepackage{booktabs}

% intermix single and multiple columns
% https://ctan.org/pkg/multicol
% \usepackage{multicol}

% create tabular cells spanning multiple rows
% https://ctan.org/pkg/multirow
\usepackage{multirow}

% extending the array and tabular environments
% https://ctan.org/pkg/array
% \usepackage{array}

% coloured boxes, for LaTeX examples and theorems, etc
% https://ctan.org/pkg/tcolorbox
\usepackage[most]{tcolorbox}

% a simple type of box for LaTeX
% https://ctan.org/pkg/minibox
\usepackage{minibox}

% variants of \fbox and other games with boxes
% https://ctan.org/pkg/fancybox
% \usepackage{fancybox}

% typeset source code listings using LaTeX
% https://ctan.org/pkg/listings
\usepackage{listings}

% place boxes at arbitrary positions on the LaTeX page
% https://ctan.org/pkg/textpos
\usepackage[absolute,overlay]{textpos}

%----------------------------------------------------------------
% cross-referencing
%----------------------------------------------------------------

% extensive support for hypertext in LaTeX
% https://ctan.org/pkg/hyperref
\usepackage{hyperref}

% Embed speaker notes in PDF metadata for WebPressive
% Notes are stored in pdfsubject using pipe-delimited format: SPEAKERNOTES|pageNum|note|pageNum|note|...
% Page numbers: 1=Title, 2=Disclaimer, 3=Outline, then slides follow in order
% NOTE: After compiling, verify page numbers match actual PDF pages and update accordingly
\hypersetup{
  pdfsubject={SPEAKERNOTES|3|This slide outlines the presentation structure, covering General, Environments, Commands, Math, Figures, Animations, and References sections.|4|This slide demonstrates a simple bulleted list layout using the myslide environment. The content uses Lorem Ipsum text to show how itemize environments work within the column structure. This is the most basic slide format in the template.|5|This slide shows a two-column layout with text on the left and a figure on the right. The myfigcol command is used to place the figure in the right column. This is a common layout for presentations where you want to combine explanatory text with visual content.|6|This slide demonstrates a vertical layout with a figure at the top and bulleted list below. The myfig command centers the figure, and myspace adds appropriate spacing between the figure and the text content. This layout is useful when the figure needs to be prominent and the text provides detailed explanation.|7|This slide shows how to place multiple figures side-by-side using the columns environment and myfigcol commands. Each figure is placed in its own column, allowing for easy comparison or showing related visual content together.|8|This slide displays all available colors in the template. The template includes general colors, grayscale options, CSU-specific colors, and Beamer theme colors that adapt to light and dark modes. All colors can be customized in the 2-colors.tex configuration file.|9|This slide demonstrates nested list environments. Both itemize (bullet points) and enumerate (numbered lists) can be nested multiple levels deep. The template automatically handles the indentation and styling for nested lists, making it easy to create hierarchical content structures.|11|This slide demonstrates the quote environment, which is used for displaying quoted text with proper indentation and italic styling. The quote environment automatically adds appropriate margins and formatting to make quoted content stand out from regular text.|12|This slide shows three special environments: mytheorem, mydefinition, and myalgorithm. These are custom environments that provide styled boxes for presenting theorems, definitions, and algorithms. Each has distinct visual styling to help differentiate types of content in technical presentations.|13|This slide demonstrates special call-out boxes that can be embedded within itemize environments using the itembox command. The template provides several box types: myremark for general remarks, myimportant for important information, myfuture for future work, and myquestion for questions. These boxes help highlight specific types of content.|14|This slide shows more advanced usage of call-out boxes, including mixing them with regular list items and using them within nested lists. The unskip command is demonstrated for better spacing control when boxes follow nested list environments. Note that this slide uses myslidefragile instead of myslide because it contains lstinline commands.|15|This slide demonstrates additional call-out box types: myoptional for optional content and myreview for review context. The myreview box accepts a parameter for the review context title. The slide also shows how these boxes can be nested within multiple levels of itemize environments.|16|This slide demonstrates code listings using the lstlisting environment. The template supports syntax highlighting for multiple languages including MATLAB and Python. The xleftmargin option is used to control indentation. Important: code listings require the myslidefragile environment instead of myslide because they contain verbatim content.|18|This slide demonstrates footnote commands: parnote for inline footnotes and parnotefull for footnotes that span the full line width. When using footnotes, you should use the 'b' option in the myslide environment to provide extra space at the bottom for the footnote content.|19|This slide demonstrates various URL commands: urlfull for complete URLs, urlhttps for HTTPS URLs without the protocol prefix, urlvideo for video links, and urlyoutube for YouTube video IDs. These commands automatically format URLs with appropriate icons and can be used in footnotes as well.|20|This slide demonstrates several utility commands: qedsymbol for the QED symbol used in proofs, myterm for highlighting important terms (which can be collected into a glossary), myline for horizontal dividers, and eqrepeat for repeating equation numbers across slides. The eqbox command is also shown for highlighting equations.|21|This slide shows various uses of the myterm command for highlighting important terms. Terms can be used inline in text, in footnotes, within theorem environments, in lists, and in call-out boxes. They can also include inline math. All terms are automatically collected and can be displayed in a glossary slide using the mytermslides command.|22|This slide demonstrates the matlabfunction command, which creates a hyperlinked reference to a MATLAB file. The command takes the filename and a URL, automatically formatting it with appropriate styling and making it clickable in the PDF.|24|This slide demonstrates the subequations environment for numbering related equations with sub-numbers (e.g., 1a, 1b, 1c). It also shows how to use notag to suppress equation numbers and tag to add custom equation labels.|25|This slide demonstrates the eqbox command, which highlights equations or mathematical expressions with a colored box. This is useful for emphasizing important equations or drawing attention to specific parts of mathematical content.|26|This slide demonstrates advanced matrix environments: bNiceMatrix for bordered matrices with row and column labels, and pNiceMatrix for parentheses matrices with various options. These environments support first-row, last-row, first-col, last-col options for labels, and special commands like Cdots and Vdots for ellipsis.|28|This slide demonstrates the myfig command, which centers a figure on the slide. The command takes the figure path and a scale factor. This is the simplest way to include a single centered figure.|29|This slide demonstrates the myfigcol command, which places figures within column environments. The first example shows a figure alongside text, while the second shows multiple figures side-by-side. The [t] option aligns columns at the top.|30|This slide demonstrates the myoverpic environment, which allows you to overlay text or graphics on top of a base figure. The put command positions elements using coordinates, and makebox with 0pt width centers text. This is useful for annotating figures with labels or additional information.|31|This slide demonstrates the myoverpiccol environment, which is similar to myoverpic but designed to work within column environments. This allows you to place annotated figures side-by-side with other content.|32|This slide demonstrates the myoverpiccolgrid environment, which is like myoverpiccol but includes a grid overlay. This is useful for aligning annotations or understanding coordinate positions when placing overlay elements.|33|This slide displays various LaTeX length dimensions used in the template, showing their values in different units (inches, millimeters, points). Understanding these dimensions is helpful for precise layout control and figure sizing.|34|This slide shows recommended figure sizes for side-by-side layouts. For general use, 2 inches width and 1.5 inches height work well. For MATLAB-generated figures, slightly larger dimensions (2.25 by 1.6875 inches) are recommended to maintain readability.|36|This slide demonstrates automatic animations using the [<+->] option in itemize environments. Each item appears sequentially when advancing slides. Note that animations only work when the handout option is removed from the documentclass. Nested lists also animate properly.|37|This slide demonstrates custom animation ordering using the uncover command with specific frame numbers. Items can appear in any order you specify (e.g., item 1, then item 4, then item 5, then item 2), and figures can also be animated. This gives precise control over the presentation sequence.|38|This slide shows that call-out boxes can also be animated using the [<+->] option. Each box and list item appears sequentially, creating a progressive reveal effect. This is useful for building up complex information step by step.|39|This slide demonstrates animating figures using the uncover command with the <+-> syntax, which automatically increments the frame number. Text and figures can appear sequentially, creating a dynamic presentation flow.|41|This slide demonstrates various BibLaTeX citation commands: cite for numeric citations, textcite for author-year in text, parencite for parenthetical citations, fullcite for complete citations, citetitle for titles, citeauthor for authors, and citeurl for URLs. Multiple citations can be combined, and citations work in footnotes. The refslides command generates a reference slide.|43|This is an appendix slide demonstrating how to include additional material after the main presentation. Appendix slides are useful for backup information, detailed proofs, or supplementary data that may be referenced during questions.}
}

% allow URL breaks at any alphanumerical character
% https://ctan.org/pkg/xurl
\usepackage{xurl}

% a new bookmark (outline) organization for hyperref
% https://ctan.org/pkg/bookmark
\usepackage{bookmark}[2019/12/03]

% line-breakable \url-like links in hyperref when compiling via dvips/ps2pdf
% https://ctan.org/pkg/breakurl
% \usepackage{breakurl}

% notes after every paragraph (or elsewhere)
% https://ctan.org/pkg/parnotes
\usepackage[restart]{parnotes}

% a range of footnote options
% https://ctan.org/pkg/footmisc
% https://tex.stackexchange.com/questions/28736
% \usepackage[para]{footmisc}

% context sensitive quotation facilities
% https://ctan.org/pkg/csquotes
\usepackage{csquotes}

%----------------------------------------------------------------
% references
%----------------------------------------------------------------

% sophisticated Bibliographies in LaTeX
% https://ctan.org/pkg/biblatex
\usepackage[style=authoryear,giveninits=true]{biblatex}

%----------------------------------------------------------------
% other
%----------------------------------------------------------------

% change format of \today with commands for current time
% https://ctan.org/pkg/datetime
\usepackage[yyyymmdd]{datetime}

% generate QR codes in LaTeX
% https://ctan.org/pkg/qrcode
\usepackage{qrcode}

% producing 'blind' text for testing
% https://ctan.org/pkg/blindtext
\usepackage{blindtext}

% 150 paragraphs of the Lorem Ipsum dummy text provided
% https://www.ctan.org/pkg/lipsum
\usepackage{lipsum}

% display various elements of a document's layout
% https://ctan.org/pkg/layouts
% \usepackage{layouts}

% create PDF and SVG animations from graphics files and inline graphics
% https://ctan.org/pkg/animate
\usepackage{animate}

% catch an external file into a macro
% https://ctan.org/pkg/catchfile
\usepackage{catchfile}